\chapter{Detailed Load Testing \& Performance Benchmarks}

\section{Introduction}
While Chapter 5 introduced the E2E testing methodology and preliminary load metrics, a production-grade system requires rigorous, multi-tiered performance profiling to ensure long-term stability under stress. This chapter expands on the system's performance characteristics, detailing the exact benchmarking suite constructed for Kemora, and providing exhaustive analytical results across varying concurrency tiers. The overarching goal is not merely to record response times, but to rigorously examine the architectural bottlenecks that limit throughput and exacerbate latency under heavy contention.

\section{Theoretical Framework: Scaling and Concurrency Limits}
Before delving into the empirical data, it is necessary to establish the theoretical boundaries that govern the scalability of the Kemora .NET backend architecture.

\subsection{Amdahl's Law and System Scalability}
Amdahl's Law posits that the theoretical maximum speedup of a system using multiple processors is limited by the strictly serial portion of the workload. Mathematically, it is expressed as:
\begin{equation}
S(N) = \frac{1}{(1-P) + \frac{P}{N}}
\end{equation}
where $S(N)$ is the theoretical speedup, $P$ is the proportion of the execution time that can be parallelized, and $N$ is the number of parallel execution units (cores or threads). In the context of the Kemora web API, the parallelizable portion ($P$) includes JSON serialization, business logic execution, and asynchronous data transformation. The serial constraints ($1-P$) encompass the initial TCP handshake, the establishment of the SSL/TLS session, Kestrel's internal request routing mechanisms, and the instantiation of the Entity Framework (EF) Core \texttt{DbContext}. As concurrency ($N$) approaches infinity, the system's throughput is ultimately bottlenecked by these sequential initialization steps, emphasizing the need for persistent connections and request pipelining.

\subsection{Thread Pool Starvation in .NET}
A critical vulnerability in highly concurrent .NET applications is Thread Pool Starvation. The Common Language Runtime (CLR) manages a global thread pool to execute asynchronous tasks. The thread pool maintains a baseline number of threads (\texttt{MinThreads}). When the volume of concurrent requests spikes, and all available threads are either executing or blocked waiting for I/O operations (such as synchronous database queries), the CLR must inject new threads. 
However, the CLR employs a pacing algorithm that restricts the thread injection rate to approximately one to two threads per second to avoid catastrophic context-switching overhead. If a sudden burst of traffic (a "thundering herd") arrives, the incoming requests are placed in a queue. If the queue time exceeds the client timeout, requests fail, even if CPU utilization remains low. The mitigation heavily relies on fully non-blocking asynchronous I/O (\texttt{async/await}) throughout the entire request pipeline, ensuring threads are yielded back to the pool while awaiting database responses.

\subsection{Connection Multiplexing and Kestrel Mechanics}
Kemora utilizes Kestrel, the cross-platform web server for ASP.NET Core. Under high concurrency, Kestrel leverages HTTP/2 connection multiplexing, allowing multiple concurrent streams to be interleaved over a single TCP connection. This drastically reduces the overhead associated with connection setup and teardown. However, connection multiplexing introduces new bottlenecks at the socket layer, specifically regarding the system's ephemeral port exhaustion and the \texttt{TIME\_WAIT} state of closed sockets. The benchmark results must be contextualized within Kestrel's bounded connection queues and the inherent limits of the underlying Windows/Linux kernel socket backlogs.

\subsection{Caching Mechanics: IMemoryCache}
To combat database latency, Kemora implements the \texttt{IMemoryCache} interface. Unlike a distributed cache (e.g., Redis), \texttt{IMemoryCache} resides within the application's memory space, providing sub-millisecond access times. Internally, it is backed by a highly optimized, lock-free \texttt{ConcurrentDictionary}. Eviction policies are governed by sliding windows and absolute expiration times, coupled with memory pressure triggers. By intercepting read-heavy requests at the application boundary, \texttt{IMemoryCache} effectively bypasses the EF Core connection pool entirely, shifting the bottleneck from database I/O to CPU-bound JSON serialization and memory bandwidth.

\section{Benchmarking Methodology}
To accurately simulate real-world traffic spikes, a custom asynchronous load-generation suite was developed in Python utilizing the \texttt{aiohttp} and \texttt{asyncio} libraries. The suite bypasses browser caching, directly stressing the Kestrel web server and the Entity Framework connection pool via REST endpoints. Concurrency is genuinely throttled: a single \texttt{asyncio.Semaphore} is instantiated once per run and acquired by every request, so the configured concurrency level reflects the true number of in-flight requests. The client also caps its TCP connection pool (\texttt{aiohttp.TCPConnector}) to the same limit, ensuring the bottleneck under test is the server rather than the load generator's own sockets.

The target endpoint is \texttt{GET /api/v1/places}, the public, anonymously-accessible places list that exercises the cache and JSON-serialization layers. Each tier ran against a freshly started, warm-cached .NET 10 server backed by SQL Server. The escalating load tiers are as follows:
\begin{itemize}
    \item \textbf{Tier 1 (Low Load):} 100 total requests at concurrency 20. Represents a sudden off-peak traffic burst.
    \item \textbf{Tier 2 (Moderate Load):} 500 total requests at concurrency 50. Represents a sustained daily active usage spike.
    \item \textbf{Tier 3 (High Load):} 1,000 total requests at concurrency 100. Represents a severe viral traffic event.
    \item \textbf{Tier 4 (Stress Test):} 2,000 total requests at concurrency 200. Designed to identify the absolute breaking point of the local connection pool.
\end{itemize}

\section{The Significance of Latency Percentiles}
In academic performance analysis, relying solely on average (mean) latency is a fundamentally flawed approach, as mean values are easily skewed by outliers and fail to represent the typical user experience. Instead, this chapter focuses on latency percentiles:
\begin{itemize}
    \item \textbf{Median (P50):} The latency experienced by 50\% of the requests. It represents the "normal" operating state of the system when cache hits are frequent.
    \item \textbf{95th Percentile (P95):} The latency experienced by the slowest 5\% of requests. It indicates the onset of queuing delays and minor resource contention.
    \item \textbf{99th Percentile (P99):} The "tail latency." The latency experienced by the slowest 1\% of requests. In distributed systems, P99 is the most critical metric. Due to "tail latency amplification," if a single user action requires fetching data from multiple microservices, the probability of experiencing the P99 latency increases exponentially. High P99 latencies are indicative of severe underlying issues such as thread pool starvation, stop-the-world garbage collection pauses, or database lock contention.
\end{itemize}

\section{Comprehensive Benchmark Results}

\subsection{Tier 1: Low Load (100 Requests, Concurrency 20)}
Under low load bursts, the system relies primarily on active memory and established SQL Server connections. There is negligible queuing delay.

\begin{table}[H]
\centering
\renewcommand{\arraystretch}{1.3}
\begin{tabular}{|l|c|p{8cm}|}
\hline
\textbf{Metric} & \textbf{Recorded Value} & \textbf{Theoretical Implication} \\ \hline
Total Requests & 100 & Baseline sample size for statistical relevance. \\ \hline
Average Latency & 47.19 ms & Low concurrency leaves the thread pool un-contended; per-request cost is dominated by JSON serialization and the cache lookup. \\ \hline
Median (P50) & 46.55 ms & Represents the raw cost of Kestrel routing + Controller execution + Cache Hit. \\ \hline
95th Percentile (P95) & 66.52 ms & Minor delays associated with asynchronous context switching. \\ \hline
99th Percentile (P99) & 67.17 ms & Near-zero spread above P95; no request paid a cache-miss penalty. \\ \hline
Error Rate & 0.00\% & System is operating well within safe boundaries. \\ \hline
\end{tabular}
\caption{Tier 1 Load Results and Theoretical Breakdown}
\end{table}

\subsection{Tier 2: Moderate Load (500 Requests, Concurrency 50)}
As the request volume increases, the system begins heavily multiplexing over the available thread pool. We observe the initial signs of resource sharing overhead.

\begin{table}[H]
\centering
\renewcommand{\arraystretch}{1.3}
\begin{tabular}{|l|c|p{8cm}|}
\hline
\textbf{Metric} & \textbf{Recorded Value} & \textbf{Theoretical Implication} \\ \hline
Total Requests & 500 & Simulates sustained daily load. \\ \hline
Average Latency & 138.18 ms & Latency rises with concurrency as requests contend for the same worker threads. \\ \hline
Median (P50) & 142.89 ms & Remains low; the cache continues to serve the read path directly. \\ \hline
95th Percentile (P95) & 213.59 ms & A controlled spread above the median indicates mild, bounded queuing rather than saturation. \\ \hline
99th Percentile (P99) & 214.78 ms & Tail latency tracks P95 closely---no pathological outliers. \\ \hline
Error Rate & 0.00\% & The CLR thread pool successfully scales up without dropping requests. \\ \hline
\end{tabular}
\caption{Tier 2 Load Results and Theoretical Breakdown}
\end{table}

\subsection{Tier 3: High Load (1,000 Requests, Concurrency 100)}
At 1,000 requests with 100-way concurrency, worker threads are saturated and incoming requests spend time queued waiting for execution, leading to a proportional rise in latency.

\begin{table}[H]
\centering
\renewcommand{\arraystretch}{1.3}
\begin{tabular}{|l|c|p{8cm}|}
\hline
\textbf{Metric} & \textbf{Recorded Value} & \textbf{Theoretical Implication} \\ \hline
Total Requests & 1,000 & Simulates a sudden, localized viral traffic spike. \\ \hline
Average Latency & 285.62 ms & Roughly doubles versus Tier 2 as queuing time accumulates under higher contention. \\ \hline
Median (P50) & 293.79 ms & The queue is now the dominant cost, not the request itself; cache hits no longer dominate the median. \\ \hline
95th Percentile (P95) & 453.66 ms & A wider but still bounded spread indicates controlled back-pressure rather than collapse. \\ \hline
99th Percentile (P99) & 456.32 ms & P99 tracks P95, confirming the absence of starvation-induced multi-second spikes. \\ \hline
Error Rate & 0.00\% & Kestrel's backlog queue prevents socket drops, maintaining 100\% reliability. \\ \hline
\end{tabular}
\caption{Tier 3 Load Results and Theoretical Breakdown}
\end{table}

\subsection{Tier 4: Stress Test (2,000 Requests, Concurrency 200)}
At 2,000 requests with 200-way concurrency, the system is subjected to a heavy stress scenario, pushing the local hardware and default CLR configurations toward their throughput ceiling.

\begin{table}[H]
\centering
\renewcommand{\arraystretch}{1.3}
\begin{tabular}{|l|c|p{8cm}|}
\hline
\textbf{Metric} & \textbf{Recorded Value} & \textbf{Theoretical Implication} \\ \hline
Total Requests & 2,000 & The upper limit exercised for the local testing environment. \\ \hline
Average Latency & 651.64 ms & Sustained queuing at 200-way concurrency; the bottleneck is thread-pool and connection-pool capacity, not the request itself. \\ \hline
Median (P50) & 678.71 ms & Latency now scales with offered load rather than staying flat---the system has entered its linear-contention regime. \\ \hline
95th Percentile (P95) & 1081.86 ms & A bounded tail; no request exceeded $\sim$1.1s, indicating graceful degradation rather than collapse. \\ \hline
99th Percentile (P99) & 1088.13 ms & P99 remains tightly coupled to P95, with no multi-second starvation spikes. \\ \hline
Error Rate & 0.00\% & Notably, \emph{zero} requests failed even at this concurrency---no timeouts, no 5xx responses, no dropped connections. \\ \hline
\end{tabular}
\caption{Tier 4 Load Results and Theoretical Breakdown}
\end{table}

\section{Analysis of the Scaling Behavior and Reliability}
The benchmark data reveals two defining characteristics of the Kemora backend: \textbf{graceful, linear latency scaling} and \textbf{perfect reliability under stress}.

\textbf{Latency scales predictably with offered load.} Across the four tiers, the median latency rises from 46.6 ms (Tier 1) to 678.7 ms (Tier 4) as concurrency increases from 20 to 200. Critically, the median tracks the average closely at every tier, and the P99 never departs far from the P95. This tight, well-behaved spread---rather than the bimodal "some instant, some multi-second" pattern that a cache-miss-dominated workload would produce---indicates that the \texttt{IMemoryCache} read path is uniformly exercised and that the dominant variable cost is thread/connection-pool queuing, not intermittent database hydration. In other words, the cache absorbs the data-access dimension, and the remaining cost is pure concurrency contention.

\textbf{Zero failures at every tier.} The most significant result is the 0.00\% error rate sustained all the way through 2,000 requests at 200-way concurrency. No request timed out, no connection was dropped, and no 5xx was returned. This demonstrates that the CLR thread pool's pacing algorithm, combined with Kestrel's bounded backlog, successfully absorbed the burst without exhausting resources---the system degrades \emph{slowly} (latency rises) rather than \emph{catastrophically} (errors spike). From a user-experience standpoint this is the preferable failure mode: under extreme load, requests get slower but still complete.

\textbf{Throughput characterization.} Treating each tier as a sustained run, the effective throughput is roughly $100/0.047 \approx 2{,}100$ req/s at Tier 1, declining as concurrency-induced queuing dominates: Tier 2 $\approx 3{,}600$ req/s of offered throughput (completed over the run), Tier 3 $\approx 3{,}500$ req/s, and Tier 4 $\approx 3{,}070$ req/s. The throughput plateaus rather than crashing, again consistent with a system bottlenecked on thread-pool injection rate---the serializable portion of Amdahl's Law---rather than on data retrieval.

\section{Hardware Utilization Observations}
During the Tier 4 stress test, local hardware utilization was monitored to identify hardware-level bottlenecks. Because the test machine was a local development host (not a dedicated, instrumented server), the following are qualitative process-level observations rather than precisely measured telemetry:
\begin{itemize}
    \item \textbf{CPU Usage:} Rose noticeably across the available logical cores during the 200-way concurrency burst. The CPU was primarily consumed by Kestrel's HTTP/2 request parsing, \texttt{System.Text.Json} serialization, and CLR garbage-collection cycles triggered by the allocation of short-lived request objects.
    \item \textbf{Memory Usage:} The API working set grew modestly from its idle baseline and remained bounded once the cache was warm, indicating effective garbage collection and the absence of a severe memory leak despite the aggressive caching policy.
    \item \textbf{Database I/O:} Disk read activity stayed low once the cache was populated, consistent with the caching strategy---the bottleneck under load was thread/connection-pool queuing and CPU cycles, not disk I/O.
\end{itemize}
For publication-grade profiling, these observations should be re-measured on a dedicated host with continuous, instrumented telemetry (e.g., \texttt{dotnet-counters} and \texttt{dotnet-trace}) to capture exact CPU\%, GC pause times, and connection-pool wait durations.

\section{Mathematical Bounding of the AI Cache Space and API Limits}
A critical factor in the system's ability to scale is the management of external AI generation requests via the OpenRouter API. Because LLM generation is inherently slow and subject to external rate limits, Kemora employs a deterministic caching strategy to bound the total number of unique AI requests required.

\subsection{Cache Permutation Mathematics}
The AI itinerary generation relies on a structured questionnaire. Assuming the system restricts the user to 3 primary questions, each with exactly 3 predefined choices (e.g., Budget: Low/Med/High; Duration: 1/3/7 Days; Style: Historic/Relax/Adventure), the maximum number of unique, distinct trip permutations ($P$) per geographic governorate is mathematically bounded:
\begin{equation}
P = C_1 \times C_2 \times C_3 = 3 \times 3 \times 3 = 27 \text{ unique combinations}
\end{equation}
If Egypt has 27 governorates ($G = 27$), the absolute theoretical maximum number of unique AI generation requests the entire system can ever experience is:
\begin{equation}
Total\_Unique\_Trips = P \times G = 27 \times 27 = 729 \text{ distinct LLM requests}
\end{equation}
Because the \texttt{OpenRouterAiService} caches the generated JSON response using a composite key of these exact parameters, the generation space is strictly mathematically bounded. By the Pigeonhole Principle, any request beyond the 729th global request is guaranteed to be a cache hit, returning an instantaneous response (0ms AI latency) instead of triggering a new 10-second OpenRouter API call. In practice, cache hits will occur much earlier due to overlapping user preferences. This bounded cache effectively neutralizes the threat of upstream AI rate-limiting at scale.

\subsection{AI Concurrency Limits and Rate Handling}
During empirical benchmarking of the AI tier itself, the system was tested against OpenRouter's upstream limits. When utilizing free-tier models, the API enforces a strict rate limit (\texttt{x-ratelimit-limit}) of approximately 20 requests per minute. Under load tests designed to bypass the cache (simulating the generation of all 729 unique permutations simultaneously), Kemora's \texttt{Polly} resilience policies successfully caught the \texttt{429 Too Many Requests} HTTP responses and applied exponential backoff with jitter. Consequently, while AI generation requests at peak load exceeded 35 seconds due to forced queuing, the application did not crash, guaranteeing eventual consistency.

\section{Conclusion on Scalability}
The benchmark suite demonstrates that the Kemora .NET 10 architecture is resilient and degrades gracefully under stress, largely due to the \texttt{IMemoryCache} implementation. The defining result is the sustained 0.00\% error rate through all four tiers---even at 2,000 requests and 200-way concurrency, every request completed successfully. Latency rose in a controlled, near-linear fashion with offered load, with tightly clustered percentiles and no multi-second starvation spikes.

For production deployment, two architectural improvements would further reduce absolute latency at high concurrency. First, implementing a distributed cache layer (such as Redis) would share the cache load and invalidate state consistently across multiple load-balanced API nodes, and would allow the read path to avoid the per-node thread-pool ceiling entirely. Second, migrating from the single-node SQL Server instance used in testing to a dedicated, heavily provisioned PostgreSQL or SQL Server cluster would increase connection-pool limits and reduce the queuing delay that currently dominates median latency at the upper tiers.
